\section{The original period vectors and their Gamma generating map}
\label{sec:OPG}

The observation in this calculation is the literal BRU1--3 matrix. We
retain its entire amplitude unit, its primary multiplicities, its
period contours, its cyclic projector and its kernel. The only source
orders used are the already present orders $s\in\{1,k\}$. The letter
$s$ in a source-order subscript is an integer parameter; the coordinate
on the one-factor polynomial algebra will be denoted $z_1$ below to
keep it distinct. The generating variable is $z$.

\subsection{The original finite algebras and the observation}
Retain
\[
\begin{gathered}
 \rho=\tfrac12+\delta+i\gamma,\quad 0<\delta<\tfrac12,\quad\gamma>2,
 \quad\mathcal R=\{\rho,\bar\rho,1-\rho,1-\bar\rho\},\\
 h(z_1)=\prod_{\alpha\in\mathcal R}(z_1-\alpha)^m,
 \quad n=4m,\quad v_h=g/h,\quad g=2\xi,\quad m\geq1,\\
 E_h=\mathbb C[z_1]/(h),\qquad A_h=M_{z_1}:E_h\to E_h,
 \qquad U=M_{j_hv_h}:E_h\to E_h.
\end{gathered}\tag{OPG1}
\]
Here $m$ is the original complete zero order at each of the four
points. Consequently $v_h(\alpha)\ne0$ for every $\alpha\in\mathcal R$.
The entire-function remainder $j_hv_h$ records its derivatives of
orders $0,\ldots,m-1$ at every one of these points. In each primary
algebra it has nonzero constant term, and its inverse is the finite
geometric series for that constant plus a nilpotent. Thus $U$ is the
same invertible unit as in BRU, with none of its values prescribed.

For the original positive integer $k$, set
\[
\begin{gathered}
 c=k/2,\quad \ell_k=1+k(m-1),\quad q=(k+1)^2\ell_k,\\
 \chi(S)=\prod_{a,b=0}^k
   \bigl(S-c-(2a-k)\delta-i(2b-k)\gamma\bigr)^{\ell_k},\\
 E=\mathbb C[S]/(\chi),\quad A=M_S:E\to E,\quad
 Y=(A-cI)/i,\\
 A_\Sigma=\sum_{j=1}^k
      I_{E_h}^{\otimes(j-1)}\otimes A_h\otimes I_{E_h}^{\otimes(k-j)},
 \qquad
 \iota:E\hookrightarrow E_h^{\otimes k},\quad
 [P]\longmapsto[P(z_1+\cdots+z_k)].
\end{gathered}\tag{OPG2}
\]
For clarity, the stated map is injective with exactly this defining
polynomial. On an ordered primary tensor the eigenvalue of
$A_\Sigma$ is the sum of the ordered roots. Its nilpotent part is
the sum of the $k$ independent nilpotents of orders $m$.
Its power $k(m-1)$ has nonzero coefficient
$(k(m-1))!/((m-1)!)^k$ on their top monomial, and its next power
is zero. The distinct sums are exactly the $(k+1)^2$ displayed
roots of $\chi$, since the signs of the real and imaginary parts
can be selected independently in every factor. Distinct pairs
$(a,b)$ give distinct sums by $\delta,\gamma>0$. Their common
nilpotency index is $\ell_k$. These facts prove that the minimal
polynomial of $A_\Sigma$ is $\chi$, and therefore prove the
kernel of polynomial evaluation and the asserted injection.
In particular
\[
 \iota A=A_\Sigma\iota,\qquad
 \iota Y=\frac{A_\Sigma-cI}{i}\iota,\qquad
 \iota[1]=[1]^{\otimes k}.
 \tag{OPG3}
\]

Keep the original $u\ne0$, its logarithm branch, and $t=0$ in
the period construction. With $\Phi_h'=h$, $\Phi_h(0)=0$,
let $\ell_j$ be the ray of angle
$(\arg u+\pi+2\pi j)/(n+1)$ and let
$\Gamma_j=-\ell_0+\ell_j$, $1\leq j\leq n$, with the original
orientations. The original period map is
\[
 \Pi:E_h\to\mathbb C^n,\qquad
 (\Pi[P])_j=\int_{\Gamma_j}e^{\Phi_h(z_1)/u}
                    \operatorname{rem}_h P(z_1)\,dz_1.
 \tag{OPG4}
\]
The representative on the right has degree less than $n$. These
are exactly the original convergent polynomial period integrals;
their existence and full determinant are proved in PDE1--13 and
IPM1--7 in the retained source closure. This definition does not
assign a period to an arbitrary representative of a quotient class.
Let $C$ be the original matrix with rows corresponding to
$\Gamma_j\mapsto\Gamma_{j+1}-\Gamma_1$ for $j<n$ and
$\Gamma_n\mapsto-\Gamma_1$. The reduced endpoint basis gives
$C^{n+1}=I$. Hence the literal finite average below is idempotent:
\[
\begin{gathered}
 P_k=\frac1{n+1}\sum_{a=0}^{n}(C^{-a})^{\otimes k},\qquad
 L=P_k\Pi^{\otimes k}U^{\otimes k}\iota,\\
 B=\operatorname{im}L,\quad\jmath:B\hookrightarrow(\mathbb C^n)^{\otimes k},
 \qquad\Lambda:E\twoheadrightarrow B,\quad\jmath\Lambda=L,
 \qquad K=\ker\Lambda.
\end{gathered}\tag{OPG5}
\]
The basis inclusion $\jmath$, the inclusion $I:K\hookrightarrow E$,
and a coefficient section $S_*:B\to E$ with $\Lambda S_*=I_B$
are the fixed BRU choices. The full actual image is used at every
rank. Empty boundary matrices use the original empty conventions.

\subsection{An exact generating identity before and after observation}
Put $b_s=s/2$, $M_s=(2\pi)^{s/2}$, and retain the monic Gamma
polynomials and their full norms from GMB1:
\[
\begin{gathered}
 p_0(y)=1,\quad p_1(y)=y,\quad
 p_{d+1}(y)=yp_d(y)-d(b_s+d-1)p_{d-1}(y),\\
 h_d^{(s)}=M_s d!(b_s)_d,\quad
 r_d=p_d(Y)[1],\quad e_d=r_d/\sqrt{h_d^{(s)}},\quad
 \lambda_d=\Lambda e_d\in B.
\end{gathered}\tag{OPG6}
\]
The square root is the positive real square root. In particular
$e_d=[u_d]$ in the original notation, including the phase $i^{-d}$
in its leading coefficient. Define the holomorphic germs of
$\log(1+z^2)$ and $\arctan z$ by their values zero at $z=0$.
For $|z|<1$, the exact algebra-valued identity is
\[
 \sum_{d=0}^{\infty}r_d\frac{z^d}{d!}
 = (1+z^2)^{-s/4}\exp((\arctan z)Y)[1].
 \tag{OPG7}
\]
Indeed the right side is holomorphic in that disc, is $[1]$ at
zero, and satisfies
$(1+z^2)F'(z)=(Y-b_s zI)F(z)$.
Comparing its Taylor coefficients gives exactly the recurrence
in OPG6: the contribution of $z^2F'$ to coefficient $z^d/d!$
is $d(d-1)r_{d-1}$, and the contribution of $b_szF$ is
$b_sd r_{d-1}$. Thus its derivatives at zero are the $r_d$,
which proves the convergent identity and not merely a formal
coefficient equality.

All factors of $A_\Sigma$ act on different tensor coordinates
and therefore commute. Their matrix exponential power series
are absolutely convergent in finite dimension. The multinomial
expansion of every power consequently proves
\[
 \exp\left(w\frac{A_\Sigma-cI}{i}\right)[1]^{\otimes k}
 =\left(\exp\left(w\frac{A_h-\tfrac12I}{i}\right)[1]\right)^{\otimes k}.
 \tag{OPG8}
\]
The centre on the left is exactly $c=k/2$; each of the $k$
centres on the right is exactly $1/2$. Applying OPG3--5 to OPG7,
and then using OPG8, gives the original observed generating map
\[
\boxed{
 \jmath\sum_{d=0}^{\infty}\sqrt{h_d^{(s)}}\lambda_d\frac{z^d}{d!}
 = (1+z^2)^{-s/4}P_k
 \left[\Pi U\exp\left((\arctan z)
                \frac{A_h-\tfrac12I}{i}\right)[1]\right]^{\otimes k}.
}\tag{OPG9}
\]
The scalar outside $P_k$ has exponent $-s/4$. It has not acquired
an extra factor of $k$: OPG7 concerns one original Gamma source
on the sum coordinate, while OPG8 is the exact factorization of
that coordinate action in the already existing tensor algebra.
For $s=1$ and $s=k$ this gives the two stated source versions of
the same original observation. No additional order is selected.

\subsection{All primary coefficients and the original period integrals}
Here is a finite formula for every vector in OPG9. For
$\alpha\in\mathcal R$ write
\[
\begin{gathered}
 h_\alpha(z_1)=h(z_1)/(z_1-\alpha)^m,\qquad
 t_\alpha(z_1)=\sum_{j=0}^{m-1}
       \frac{(h_\alpha^{-1})^{(j)}(\alpha)}{j!}(z_1-\alpha)^j,\\
 \varepsilon_\alpha=[h_\alpha t_\alpha]\in E_h,
 \qquad E_{\alpha,a}(z_1)=
   \operatorname{rem}_h\bigl(h_\alpha(z_1)t_\alpha(z_1)
                                  (z_1-\alpha)^a\bigr),\quad0\leq a<m,\\
 \mathbf b_{\alpha,a}=\Pi[E_{\alpha,a}],\qquad
 (\mathbf b_{\alpha,a})_j=
   \int_{\Gamma_j}e^{\Phi_h(z_1)/u}E_{\alpha,a}(z_1)\,dz_1.
\end{gathered}\tag{OPG10}
\]
The Taylor inverse exists because $h_\alpha(\alpha)\ne0$.
Reduction modulo $(z_1-\alpha)^m$ gives
$\varepsilon_\alpha=1$ on its primary factor and zero on every
other primary factor. These observations prove
$\varepsilon_\alpha\varepsilon_\beta=0$ for $\alpha\ne\beta$,
$\varepsilon_\alpha^2=\varepsilon_\alpha$, and
$\sum_\alpha\varepsilon_\alpha=1$: their differences are divisible
by all four coprime factors and hence by $h$. They also prove that
$[E_{\alpha,a}]=\varepsilon_\alpha(z_1-\alpha)^a$ is a basis of
$E_h$. Every representative in OPG10 has degree less than $n$,
so each integral is precisely of the original polynomial type.

For $w\in\mathbb C$ put
$\mathbf v(w)=\Pi U\exp(w(A_h-\tfrac12I)/i)[1]$.
The primary nilpotent $N_\alpha=(z_1-\alpha)\varepsilon_\alpha$
satisfies $N_\alpha^m=0$. Taylor multiplication in that algebra
therefore proves, with all amplitude derivatives retained,
\[
\boxed{
 \mathbf v(w)=\sum_{\alpha\in\mathcal R}
 e^{w(\alpha-1/2)/i}\sum_{a=0}^{m-1}\mathbf b_{\alpha,a}
       \sum_{\nu=0}^{a}
       \frac{v_h^{(\nu)}(\alpha)}{\nu!}
       \frac{(w/i)^{a-\nu}}{(a-\nu)!}.
}\tag{OPG11}
\]
For a direct verification before integration, both
$U\exp(w(A_h-\tfrac12I)/i)[1]$ and
$j_h(v_h(z_1)e^{w(z_1-1/2)/i})$ have the same first $m$ Taylor
coefficients at every $\alpha$. Their difference is therefore
zero in $E_h$. The representative to which $\Pi$ is applied is
exactly the finite polynomial in the basis $E_{\alpha,a}$ whose
coefficients are displayed in OPG11. In particular this proof
does not replace the integrand in OPG10 by the unreduced entire
function $v_h(z_1)e^{w(z_1-1/2)/i}$. The precise map from that
entire function to the period vector is its finite Taylor
remainder $j_h$, followed by the degree-less-than-$n$ representative,
followed by the original polynomial period integral.

The cyclic projector is also fully evaluated as a finite operation:
\[
 \jmath\sum_{d\geq0}\sqrt{h_d^{(s)}}\lambda_d\frac{z^d}{d!}
 =\frac{(1+z^2)^{-s/4}}{n+1}
       \sum_{a=0}^{n}\bigl(C^{-a}\mathbf v(\arctan z)\bigr)^{\otimes k}.
 \tag{OPG12}
\]
This is obtained by applying each tensor power in OPG5 to a pure
tensor. It retains every term of the original average, its factor
$1/(n+1)$, and its inverse powers of $C$.

There is also a finite coefficient formula requiring no series
extraction. For an ordered $k$-tuple
$\boldsymbol\alpha\in\mathcal R^k$ set
$\omega_{\boldsymbol\alpha}=(\sum_i\alpha_i-c)/i$.
For $0\leq a_i<m$ and $0\leq\nu_i\leq a_i$ put
$D(\mathbf a,\boldsymbol\nu)=\sum_i(a_i-\nu_i)$ and
\[
 T_d(\boldsymbol\alpha,\mathbf a)=
 \sum_{\substack{0\leq\nu_i\leq a_i\\1\leq i\leq k}}
 \frac{i^{-D(\mathbf a,\boldsymbol\nu)}
        p_d^{(D(\mathbf a,\boldsymbol\nu))}(\omega_{\boldsymbol\alpha})}
        {\prod_{i=1}^k(a_i-\nu_i)!}
 \prod_{i=1}^k\frac{v_h^{(\nu_i)}(\alpha_i)}{\nu_i!}.
 \tag{OPG13}
\]
As usual the polynomial derivative is zero when its order exceeds
$d$. Multivariate Taylor expansion in the independent nilpotents
$(z_i-\alpha_i)$ proves the exact formula
\[
\boxed{
 \jmath\lambda_d=
 \frac1{\sqrt{h_d^{(s)}}}\frac1{n+1}
 \sum_{a=0}^{n}\ \sum_{\boldsymbol\alpha\in\mathcal R^k}
 \sum_{\substack{0\leq a_i<m\\1\leq i\leq k}}
 T_d(\boldsymbol\alpha,\mathbf a)
 \bigotimes_{i=1}^k C^{-a}\mathbf b_{\alpha_i,a_i}.
}\tag{OPG14}
\]
Indeed on that primary tensor, expand each $v_h(z_i)$ first.
For the remaining multi-degree $a_i-\nu_i$, Taylor expansion of
$p_d((\sum_i z_i-c)/i)$ supplies its derivative of the total
order $D$, the factor $i^{-D}$, and the factorials
$\prod_i(a_i-\nu_i)!$. These are exactly the coefficients in
OPG13. Apply the $k$ polynomial period maps and the literal cyclic
average to get OPG14. All sums are finite. Ordered tuples have
not been replaced by just their eigenvalue sum, so both the
amplitude data and every independent primary nilpotent survive.

\subsection{The actual observation increment without a coordinate inverse}
Retain the GMB maps
$C_s=[e_0,\ldots,e_{q-1}]:\mathbb C^q\to E$,
$v_j=C_s^{-1}e_{q+j}$,
$K_j=I_q+\sum_{a<j}v_av_a^*$, and $R_s=\Lambda C_s$.
The low-degree leading coefficients prove $C_s$ invertible as in
GMB2. On each column, applying the maps gives
\[
 R_s v_j=\lambda_{q+j},\qquad
 R_sK_jR_s^*=\sum_{d=0}^{q+j-1}\lambda_d\lambda_d^*.
 \tag{OPG15}
\]
For the second equality the identity term in $K_j$ gives exactly
the first $q$ columns of $R_s$, and each remaining rank-one term
gives the next displayed column. The low vectors $\lambda_d$,
$0\leq d<q$, span $B$ because $C_s$ is invertible and $\Lambda$
is onto. Thus, in the fixed original basis of $B$, the matrix
\[
 \mathcal M_D=\sum_{d=0}^{D-1}\lambda_d\lambda_d^*\quad(D\geq q)
 \tag{OPG16}
\]
is positive definite whenever $B$ has positive dimension:
for a nonzero covector its quadratic form is the sum of its
squared values on spanning vectors. Consequently the exact
original boundary increment of GMB8 is
\[
 \boxed{\displaystyle
 \xi_j=\lambda_{q+j}^*\mathcal M_{q+j}^{-1}\lambda_{q+j},\qquad
 1+\xi_j=\frac{\det\mathcal M_{q+j+1}}{\det\mathcal M_{q+j}}.}
 \tag{OPG17}
\]
The determinant equality follows by conjugating the rank-one
update by $\mathcal M_{q+j}^{-1/2}$; its sole possible nonzero
added eigenvalue is the displayed quadratic form. If $B=0$,
the vector is empty, $\xi_j=0$ and both determinants equal one.
Thus OPG11--14 compute the boundary increment from the original
period vectors, with no $C_s^{-1}$ in OPG17. This cancellation
is an identity of the proved maps, rather than a change of the
observation or the source measure.

At the four original endpoints the boundary contribution is
therefore the exact finite expression
\[
 F_B^{(s)}=
 \log\frac{\det\mathcal M_{2q}\det\mathcal M_{2q+1}}
                 {\det\mathcal M_q\det\mathcal M_{q+1}}
 =\log(1+\xi_0)+2\sum_{j=1}^{q-1}\log(1+\xi_j)
                         +\log(1+\xi_q).
 \tag{OPG18}
\]
Expanding the last sum of consecutive log determinant differences
leaves coefficient $-1$ at $\mathcal M_q,\mathcal M_{q+1}$ and
$+1$ at $\mathcal M_{2q},\mathcal M_{2q+1}$, proving every sign.
GMB13 proves that the other original factor is
$(1+t_j)/(1+\xi_j)$, with $t_j=v_j^*K_j^{-1}v_j$ and
$0\leq\xi_j\leq t_j$. Thus OPG18 is the actual BRU boundary
contribution, and the full original kernel contribution remains
$\mathcal W_q(\log((1+t_j)/(1+\xi_j)))$.

\subsection{A finite recurrence retaining the exact action defect}
The period formula also has a finite recurrence on the fixed
original coefficient spaces. Define the uniquely typed projection
$\kappa:E\to K$ by
$I\kappa=I_E-S_*\Lambda$. It exists because the right side has
image in $\ker\Lambda=\operatorname{im}I$, and $I$ is an
inclusion. It satisfies $\kappa I=I_K$ and $\kappa S_*=0$.
In particular $E=I(K)\oplus S_*(B)$ as coefficient spaces, without
an assertion that this splitting is preserved by $Y$.
Set
\[
\begin{gathered}
 \mathsf Y_B=\Lambda YS_*:B\to B,\quad
 \mathsf D=\Lambda YI:K\to B,\\
 \mathsf C=\kappa YS_*:B\to K,\quad
 \mathsf Y_K=\kappa YI:K\to K,\quad
 \zeta_d=\kappa e_d\in K.
\end{gathered}\tag{OPG19}
\]
All four maps are computable from OPG1--5 and the fixed BRU
section. For $d\geq0$ put
\[
 a_d=\frac1{\sqrt{(d+1)(b_s+d)}},\qquad
 b_d=\sqrt{\frac{d(b_s+d-1)}{(d+1)(b_s+d)}}\quad(d\geq1),
 \qquad b_0=0.
 \tag{OPG20}
\]
Set $\lambda_{-1}=0$ and $\zeta_{-1}=0$. Then the exact coupled
recurrence is
\[
\boxed{\begin{aligned}
 \lambda_{d+1}&=a_d(\mathsf Y_B\lambda_d+\mathsf D\zeta_d)
                                          -b_d\lambda_{d-1},\\
 \zeta_{d+1}&=a_d(\mathsf C\lambda_d+\mathsf Y_K\zeta_d)
                                          -b_d\zeta_{d-1},\\
 \lambda_0&=\Lambda[1]/\sqrt{M_s},\qquad
 \zeta_0=\kappa[1]/\sqrt{M_s}.
\end{aligned}}\tag{OPG21}
\]
To prove this, the norm ratio from OPG6 gives
$\sqrt{h_d^{(s)}/h_{d+1}^{(s)}}=a_d$ and
$d(b_s+d-1)\sqrt{h_{d-1}^{(s)}/h_{d+1}^{(s)}}=b_d$.
Hence $e_{d+1}=a_dYe_d-b_de_{d-1}$, including $d=0$ by its
initial values. Substituting $e_d=S_*\lambda_d+I\zeta_d$ and
applying $\Lambda$ and $\kappa$ proves both lines of OPG21.
It retains the kernel coordinate at every degree and supplies
the period vectors through degree $2q$, exactly the range used
in OPG17--18.

The failure of a putative boundary action has the precise map
\[
 \Lambda Y-\mathsf Y_B\Lambda=\mathsf D\kappa:E\to B,\qquad
 \jmath\mathsf D=
 \frac1iP_k\Pi^{\otimes k}U^{\otimes k}
                   (A_\Sigma-cI)\iota I:K\to(\mathbb C^n)^{\otimes k}.
 \tag{OPG22}
\]
The first equality follows from $I_E=S_*\Lambda+I\kappa$;
the second follows from OPG3--5. These equations preserve the
cyclic average and the full unit. They also prove exactly that
$Y(K)\subseteq K$ holds if and only if $\mathsf D=0$.
No such vanishing is used anywhere in OPG7--21. Thus every
possible boundary action defect is retained inside a proved
coupled action and inside the explicit original period formula.

For completeness the multiplication action does descend to a
canonical quotient that remembers every observed iterate. Define
\[
 \mathcal K_\infty=\bigcap_{a=0}^{q-1}\ker(\Lambda Y^a),\qquad
 \mathcal O:E\to B^{\oplus q},\quad
 x\longmapsto(\Lambda x,\Lambda Yx,\ldots,\Lambda Y^{q-1}x).
 \tag{OPG23}
\]
The polynomial $\chi(c+iT)$ annihilates $Y$ by OPG2, has degree
$q$ and leading coefficient $i^q\ne0$. It therefore expresses
$Y^q$ as a linear combination of $I,Y,\ldots,Y^{q-1}$, with
the exact coefficients of this original polynomial. If
$x\in\mathcal K_\infty$, each component of $\mathcal O(Yx)$
vanishes by this relation. Thus $\mathcal K_\infty$ is
$Y$-invariant and is precisely $\bigcap_{a\geq0}\ker(\Lambda Y^a)$.
Moreover any $Y$-invariant subspace contained in $K$ lies in
this intersection. Consequently the first isomorphism map
\[
\begin{gathered}
 E/\mathcal K_\infty\ \xrightarrow{\ \overline{\mathcal O}\ }
 \operatorname{im}\mathcal O,\quad[x]\longmapsto\mathcal O(x),\\
 \chi_j=[T^j]\chi(c+iT),\quad\chi_q=i^q,\\
 \mathsf{Shift}(b_0,\ldots,b_{q-1})=
 \left(b_1,\ldots,b_{q-1},-i^{-q}\sum_{j=0}^{q-1}\chi_j b_j\right).
\end{gathered}
 \tag{OPG24}
\]
is a bijection intertwining multiplication by $Y$ with the shift
of the $q$ components, whose last component is the stated exact
linear combination. It projects to the original observation by
the first component, with retained kernel $K/\mathcal K_\infty$.
This proves an action-preserving enlargement of the observed
data on the original algebra. It neither discards the original
kernel in OPG21 nor presumes a vanished defect in OPG22.

The original period isomorphism also locates that defect in the
cyclic projection. Its inverse exists by the retained period
determinant theorem used in OPG4. On the unchanged ambient tensor
space put
\[
 \mathcal V=\Pi^{\otimes k}U^{\otimes k}\iota,
 \qquad
 \mathcal T=\sum_{j=1}^k
 \left(\Pi\frac{A_h-\tfrac12I}{i}\Pi^{-1}\right)_j.
 \tag{OPG25}
\]
The subscript $j$ means that the displayed factor acts in position
$j$ and identities act in all other positions. Multiplication by
$j_hv_h$ commutes with $A_h$. Equations OPG3 and OPG25 therefore
give $\mathcal VY=\mathcal T\mathcal V$ exactly. Since
$L=P_k\mathcal V$ and $P_k^2=P_k$, splitting the ambient identity
as $P_k+(I-P_k)$ now proves
\[
\begin{gathered}
 LY=P_k\mathcal T L+P_k\mathcal T(I-P_k)\mathcal V,\\
 \jmath\mathsf D=P_k\mathcal T(I-P_k)\mathcal V I.
\end{gathered}\tag{OPG26}
\]
For the second equality, $P_k\mathcal V I=LI=0$, so the
first term of the first equality vanishes on $I(K)$. This is
the same defect as OPG22 on its original domain. The two ambient
terms in the first line sum to a vector in $\jmath(B)$; neither
individual term is asserted to preserve that image. OPG19--21
give the full recurrence within $K\oplus B$ without such an
assertion. Thus the preserved unit and period isomorphism carry
the multiplication action exactly, and OPG26 identifies precisely
the retained contribution of the original cyclic projection.

Equations OPG9, OPG11--14, OPG17 and OPG21 now give the same
original boundary increments by analytic generating functions,
finite primary coefficients, and finite coupled recurrences.
Their derivations establish the exact maps connecting these
calculations. In particular the new input into the original
mixed-control problem is the fully specified sequence
$\lambda_0,\ldots,\lambda_{2q}$ with its actual unit and periods,
rather than a freely selected observation or a source-order
replacement. The formulas themselves do not assign a growing-
parameter estimate to the individual weighted sums; they supply
their unchanged mathematical input for that calculation.

% END RX INLINE 16
